\documentclass[a4paper]{article}

%Tutti gli usepackage vanno qui

\usepackage{geometry}
\usepackage[italian]{babel}
\usepackage[utf8]{inputenc}
\usepackage[T1]{fontenc}
\usepackage{tgschola}
%\usepackage{tgbonum}
\usepackage{tabularx}
\usepackage{longtable}
\usepackage{hyperref}
\usepackage{enumitem}
\usepackage[toc]{appendix}
\hypersetup{
	colorlinks=true,
	linkcolor=black,
	filecolor=magenta,      
	urlcolor=blue,
}
% Numerazione figure 
\let\counterwithout\relax
\let\counterwithin\relax
\usepackage{chngcntr}

\counterwithin{table}{subsection}
\counterwithin{figure}{subsection}

\usepackage[bottom]{footmisc}
\usepackage{fancyhdr}
\setcounter{secnumdepth}{4}
\usepackage{amsmath, amssymb}
\usepackage{array}
\usepackage{graphicx}

\usepackage{ifthen}

%\usepackage{float}
\usepackage{layouts}
\usepackage{url}
\usepackage{comment}
\usepackage{float}
\usepackage{eurosym}

\usepackage{lastpage}
\usepackage{layouts}
\usepackage{float}
\usepackage{eurosym}

%Comandi di impaginazione uguale per tutti i documenti
\pagestyle{fancy}
\lhead{\includegraphics[scale=0.04]{../../../../latex/images/logoTWM.png}}
%Titolo del documento
\rhead{\doctitle{}}
%\rfoot{\thepage}
\cfoot{Pagina \thepage\ di \pageref{LastPage}}
\setlength{\headheight}{35pt}
\setcounter{tocdepth}{5}
\setcounter{secnumdepth}{5}
\renewcommand{\footrulewidth}{0.4pt}

% multirow per tabelle
\usepackage{multirow}

% Permette tabelle su più pagine
%\usepackage{longtable}


% colore di sfondo per le celle
\usepackage[table]{xcolor}

%COMANDI TABELLE
\newcommand{\rowcolorhead}{\rowcolor[HTML]{9b240a}} %intestazione 
% check for missing commands
\newcommand{\headertitle}[1]{\textbf{\color{white}#1}} %titolo colonna
\definecolor{pari}{HTML}{FFDBCB}
\definecolor{dispari}{HTML}{F1F7FD}

% comandi glossario
\newcommand{\glo}{$_{G}$}
\newcommand{\glosp}{$_{G}$ }


%label custom
\makeatletter
\newcommand{\uclabel}[2]{%
	\protected@write \@auxout {}{\string \newlabel {#1}{{#2}{\thepage}{#2}{#1}{}} }%
	\hypertarget{#1}{#2}
}
\makeatother

%riportare pezzi di codice
\definecolor{codegray}{gray}{0.9}
\newcommand{\code}[1]{\colorbox{codegray}{\texttt{#1}}}



% Configurazione della pagina iniziale
\newcommand{\doctitle}{Norme di Progetto}
\newcommand{\docdate}{ } % lasciare vuoto (uno carattere di spazio) per i documenti che non sono verbali, così non viene scritta la data
\newcommand{\rev}{3.0.1}
\newcommand{\stato}{Non Approvato}
\newcommand{\uso}{Interno}
\newcommand{\approv}{--De Renzis Simone}
\newcommand{\red}{Greggio Nicolò\\& Tessari Andrea}
\newcommand{\ver}{Zuccolo Giada\\& Crivellari Alberto}
\newcommand{\dest}{Three Way Milkshake\\ Prof. Vardanega Tullio\\ Prof.
 Cardin Riccardo}
\newcommand{\describedoc}{Questo documento contiene tutte le norme di progetto\textsubscript{G}, definite inizialmente o aggiunte in seguito, viene quindi aggiornato per incrementi successivi.}



 % modifica questo file
\makeindex

\usepackage{hyperref}
\hypersetup{
    colorlinks=true,
    linkcolor=blue,
    urlcolor=blue,
    hyperfootnotes=false
}
\usepackage{multicol}
\usepackage{pgfplots}
\usepackage{verbatim}
\usepackage{pgf-pie}
\usepackage{ragged2e}
\setlength{\columnseprule}{1pt}


\begin{document}
	\thispagestyle{empty}
\begin{titlepage}
	\begin{center}
		
		\includegraphics[scale = 0.17]{../../../../latex/images/logoTWM.png}\\[0.7cm]
		

		\noindent\rule{\textwidth}{1pt} \\[0.4cm]
		\Huge \textbf{\doctitle} \\[0.1cm]
		\ifthenelse{\equal{\docdate}{ }}{ }{ \huge \textbf{\docdate} \\[0.1cm] }
		
		\noindent\rule{\textwidth}{1pt}\\[0.7cm]
		
		\large \textbf{Three Way Milkshake - Progetto "PORTACS"} \\[0.4cm] 
                \texttt{threewaymilkshake@gmail.com} \\[0.4cm]
                
		
        
        
        \large

        \begin{tabular}{r|l}
                        \textbf{Versione} & \rev{} \\
                        \textbf{Stato} & \stato{} \\
                        \textbf{Uso} & \uso{} \\                         
                        \textbf{Approvazione} & \approv{} \\                      
                        \textbf{Redazione} & \red{} \\ 
                        \textbf{Verifica} &  \ver{} \\                         
                        \textbf{Destinatari} & \parbox[t]{5cm}{ \dest{} }
                \end{tabular} 
                \\[0.3cm]
                \large \textbf{Descrizione} \\ \describedoc{} 
               

	\end{center}
\end{titlepage}
	\pagebreak

	% Registro delle modifiche
	\section*{Registro delle modifiche}

\newcommand{\changelogTable}[1]{
	
	
	\renewcommand{\arraystretch}{1.5}
	\rowcolors{2}{pari}{dispari}
	\begin{longtable}{ 
			>{\centering}p{0.07\textwidth} 
			>{}p{0.21\textwidth}
			>{\centering}p{0.17\textwidth}
			>{\centering}p{0.13\textwidth} 
			>{\centering}p{0.17\textwidth} 
			>{\centering}p{0.13\textwidth} }
		\rowcolorhead
		\headertitle{Vers.} &
		\centering \headertitle{Descrizione} &	
		\headertitle{Redazione} &
		\headertitle{Data red.} & 
		\headertitle{Verifica} &
		\headertitle{Data ver.}
		\endfirsthead	
		\endhead
		
		#1
		
	\end{longtable}
	\vspace{-2em}
	
}


\newcommand{\approvingTable}[1]{ 
	
	
	\renewcommand{\arraystretch}{1.5}
	\rowcolors{2}{pari}{dispari}
	\begin{longtable}{ 
			>{\centering}p{0.07\textwidth} 
			>{\centering}p{0.415\textwidth}
			>{\centering}p{0.13\textwidth}
			>{\centering}p{0.322\textwidth}  }
		\rowcolorhead
		\headertitle{Vers.} &
		\centering \headertitle{Descrizione} &	
		\headertitle{Data appr.} &
		\headertitle{Approvazione}
		\endfirsthead	
		\endhead
		
		#1
		
	\end{longtable}
	\vspace{-2em}
	
}
	\changelogTable{
	3.5.0 & Incremento appendice \S\ A e B & Crivellari Alberto & 2021-04-26 & -- & --
	3.4.0 & Rimossa appendice \S\ D & Crivellari Alberto & 2021-04-01 & -- & --
	\tabularnewline
	3.3.0 & Modifica sezioni \S\ 2 e \S\ 3 & Crivellari Alberto & 2021-04-01 & -- & --
	\tabularnewline
	3.2.0 & Modifica appendice \S\ B & Zuccolo Giada & 2021-03-31 & -- & --
\tabularnewline
	3.1.0 & Aggiornamento appendice \S\ C & Crivellari Alberto & 2021-03-20 & -- & --
}

\approvingTable{
    3.0.0 & Approvazione del documento & 2021-03-08 & Zuccolo Giada
}

\changelogTable{
    2.1.0 & Aggiornamento sezioni \S\ 4 e 5 & Crivellari Alberto & 2021-03-08 & Chiarello Sofia & 2021-03-08
}

\approvingTable{
    2.0.0 & Approvazione del documento & 2021-02-22 & De Renzis Simone
}

\changelogTable{
1.3.0 & Completamento stesura appendice \S\ D e \S\ B & Crivellari Alberto & 2021-02-12 & Greggio Nicolò & 2021-02-19
\tabularnewline
1.2.1 & Completamento stesura sezione \S\ 2 e \S\ 3 & Chiarello Sofia & 2021-02-11 & Tessari Andrea & 2021-02-13
\tabularnewline
1.2.0 & Stesura sezione \S\ 3 e appendice \S\ C & Chiarello Sofia & 2021-02-09 & Tessari Andrea & 2021-02-12
\tabularnewline
1.1.0 & Stesura sezione \S\ 2 & Chiarello Sofia & 2021-02-08 & Tessari Andrea & 2021-02-13
}

\approvingTable{
	1.0.0 & Approvazione del documento & 2021-01-10 & De Renzis Simone
}

\changelogTable{
0.5.1 & Aggiunta tabelle a sezione \S 4 e \S5 & Crivellari Alberto & 2021-01-08 & Greggio Nicolò & 2021-01-10
\tabularnewline
0.5.0 & Redazione sezione \S 5 & Crivellari Alberto & 2020-12-07 & Greggio Nicolò & 2021-01-10
\tabularnewline
0.4.0 & Redazione sezione \S 4 & Crivellari Alberto & 2020-12-06 & Greggio Nicolò & 2021-01-10
\tabularnewline
0.3.2 & Modifiche sezione \S 1 & Crivellari Alberto & 2020-12-30 & Greggio Nicolò & 2021-01-10
\tabularnewline
0.3.1 & Tabelle sezione \S 3 & Crivellari Alberto & 2020-12-29 & Greggio Nicolò & 2021-01-10
\tabularnewline
0.3.0 & Redazione sezione \S 3 & Tessari Andrea & 2020-12-28 & Greggio Nicolò & 2021-01-10
\tabularnewline
0.2.1 & Tabelle sezione \S 2 & Crivellari Alberto & 2020-12-20 &  Greggio Nicolò & 2021-01-09
\tabularnewline
0.2.0 & Redazione sezione \S 2 & Crivellari Alberto & 2020-12-19 & Greggio Nicolò & 2021-01-09
\tabularnewline
0.1.1 & Modifiche sezione \S 1  & Crivellari Alberto & 2020-12-18 & Greggio Nicolò & 2021-01-09
\tabularnewline
0.1.0 & Redazione sezione \S 1 & Crivellari Alberto & 2020-12-16 & Greggio Nicolò & 2021-01-09
\tabularnewline
0.0.1 & Strutturazione del documento & Tessari Andrea & 2020-12-15 & Greggio Nicolò & 2021-01-09
}


 % modifica questo file
	%\end{longtable}
	\pagebreak

	% indice
	{
        \hypersetup{linkcolor=black}
        \tableofcontents
        \pagebreak

        % indice delle figure
        \listoffigures
        \pagebreak

        % indice delle tabelle
        \listoftables
        \pagebreak
    }

	\newcommand{\contabilitaTable}[1]{

\begin{table}[H]
	\begin{center}
		\begin{tabular}{c
				!{\color[HTML]{9b240a}\vrule width 1pt}
				cccccc
				!{\color[HTML]{9b240a}\vrule width 1pt}	
				c}
			\rowcolorhead
			\headertitle{Nome} & \headertitle{R} & \headertitle{V} & \headertitle{An} & \headertitle{Am} & \headertitle{Pr} & \headertitle{Pt} & \headertitle{Tot} \\	
			#1
			\end{center}
			\end{table}	
}


\newcommand{\smallPreventivoTable}[1]{
	
	\begin{table}[H]
		\begin{center}
			\begin{tabular}{c
					!{\color[HTML]{9b240a}\vrule width 1pt}
					cccccc
					!{\color[HTML]{9b240a}\vrule width 1pt}	
					c}
				\rowcolorhead
				\headertitle{Ruolo} & \headertitle{Tempo (ore)} & \headertitle{Costo (\euro)} \\
				#1
			\end{center}
		\end{table}	
	}


\newcommand{\planningTable}[1]{
	
	
	\renewcommand{\arraystretch}{1.5}
	\rowcolors{2}{pari}{dispari}
	\begin{longtable}{ 
			>{}p{0.25\textwidth} 
			>{}p{0.42\textwidth}
			>{\centering}p{0.05\textwidth}
			>{\centering}p{0.17\textwidth} }
		\rowcolorhead
		\headertitle{Attività} &
		\headertitle{Descrizione} &	
		\headertitle{Ore} &
		\headertitle{Ruolo} 
		\endfirsthead	
		\endhead
		#1		
	\end{longtable}
	
}


\newcommand{\pafTable}[1]{
	
	
	\renewcommand{\arraystretch}{1.5}
	\rowcolors{2}{pari}{dispari}
	\begin{longtable}{ 
			>{}p{0.40\textwidth} 
			>{\centering}p{0.10\textwidth}
			>{\centering}p{0.20\textwidth}
			>{\centering}p{0.15\textwidth} }
		\rowcolorhead
		\headertitle{Fase} &
		\headertitle{Periodo} &	
		\headertitle{Descrizione} &
		\headertitle{Spesa (\euro)} 
		\endfirsthead	
		\endhead
		#1		
	\end{longtable}
	
}


 % file con template tabelle con macro
	

	% contenuto del documento, ogni sezione in un file
	\section{Introduzione}
\subsection{Scopo del documento}
    Questo documento ha lo scopo di fissare e definire tutte le regole, convenzioni e buone pratiche utili a formare un way of working condiviso alla base da tutti i componenti del gruppo per assicurare una collaborazione efficiente ed efficace. Si discuteranno inoltre i vari strumenti che verranno adottati per facilitare lo sviluppo del progetto\textsubscript{G} e per promuovere un'organizzazione adeguata.
    Si ritiene inoltre che la definizione ed il mantenimento per incremento di un documento condiviso all'interno del gruppo di lavoro, che definisca e raccolga quanto descritto in maniera formale e centralizzata, possa favorire, in un contesto dove i membri possano variare, l'inserimento di nuovi componenti facilitandone l'ambientamento. Pur non essendo questo il contesto di lavoro attuale, è comunque una buona pratica da sperimentare e consolidare.

\subsection{Scopo del prodotto}
Il capitolato\textsubscript{G} C5 propone un progetto\textsubscript{G} in cui viene richiesto lo sviluppo di un software per il monitoraggio in tempo reale di unità che si muovono in uno spazio definito. All'interno di questo spazio, creato dall'utente per riprodurre le caratteristiche di un ambiente reale, le unità dovranno essere in grado di circolare in autonomia, o sotto il controllo dell'utente, per raggiungere dei punti di interesse posti nella mappa.  La circolazione è sottoposta a vincoli di viabilità e ad ostacoli propri della topologia dell'ambiente, deve evitare le collisioni con le altre unità e prevedere la gestione di situazioni critiche nel traffico.

\subsection{Termini, abbreviazioni ed altri documenti}
    Tutti i termini che necessitano di una spiegazione, per fornire un'adeguata comprensione, o perché possono causare ambiguità nel contesto, sono definiti nel glossario. A dispetto di tutti gli altri documenti presenti, il glossario è sotto la veste di una pagina web, più nello specifico in una wiki\textsubscript{G} di GitHub, accessibile a \href{https://github.com/Three-Way-Milkshake/docs/wiki/Glossario}{questo indirizzo}, con vocaboli suddivisi tra "Acronimi" e "Termini", disposti in ordine alfabetico al fine di facilitare la navigazione. All'interno dei documenti le voci di glossario saranno seguite da una G pedice mentre gli acronimi da una A (e.g.: voce di glossario\textsubscript{G} ; acronimo\textsubscript{A}).
    Inoltre quando si farà riferimento ad un altro documento o al documento stesso, il nome di questo sarà in maiuscoletto (e.g.: \textsc{esempio nome documento}).

\subsection{Riferimenti}
\label{ref}
    \subsubsection{Riferimenti Normativi}
        \begin{itemize}
            \item Standard ISO 12207:\\ \url{https://www.math.unipd.it/~tullio/IS-1/2009/Approfondimenti/ISO_12207-1995.pdf};
            \item Standard UML\textsubscript{A} 2.0:\\\url{https://www.omg.org/spec/UML/2.0/Superstructure/PDF};
            \item Diagrammi dei casi d'uso\textsubscript{G}:\\\url{https://www.math.unipd.it/~rcardin/swea/2021/Diagrammi\%20Use\%20Case_4x4.pdf};
            \item Diagrammi delle classi:\\\url{https://www.math.unipd.it/~rcardin/swea/2021/Diagrammi\%20delle\%20Classi_4x4.pdf};
            \item Diagrammi dei package:\\\url{https://www.math.unipd.it/~rcardin/swea/2021/Diagrammi\%20dei\%20Package_4x4.pdf};
            \item Diagrammi di attività\textsubscript{G}:\\\url{https://www.math.unipd.it/~rcardin/swea/2021/Diagrammi\%20di\%20Attivit\%c3\%a0_4x4.pdf};
            \item Diagrammi di sequenza:\\\url{https://www.math.unipd.it/~rcardin/swea/2021/Diagrammi\%20di\%20Sequenza_4x4.pdf};
            \item Standard ISO 8601:\\\url{https://www.iso.org/iso-8601-date-and-time-format.html}.
        \end{itemize}

    \subsubsection{Riferimenti Informativi}
        \begin{itemize}
        	\item \textsc{\href{https://github.com/Three-Way-Milkshake/docs/wiki/Glossario}{Glossario}}:\\per la definizione dei termini (pedice G) e degli acronimi (pedice A) evidenziati nel documento;
            \item \url{https://www.math.unipd.it/~tullio/IS-1/2020/Dispense/L03.pdf}
            \item \url{https://www.math.unipd.it/~tullio/IS-1/2020/Dispense/L06.pdf}
            \item \url{https://www.math.unipd.it/~tullio/IS-1/2020/Dispense/FC2.pdf}
            \item \url{https://www.math.unipd.it/~rcardin/swea/2021/SOLID\%20Principles\%20of\%20Object-Oriented\%20Design_4x4.pdf}
            \item \url{https://www.math.unipd.it/~tullio/IS-1/2020/Dispense/L09.pdf}
            \item \url{https://docs.github.com/en/free-pro-team@latest/actions}
            \item \url{https://www.latex-project.org/}
            \item \url{http://www.texstudio.org/}
            \item \url{https://www.xm1math.net/texmaker/}
            \item \url{https://www.atlassian.com/git/tutorials/comparing-workflows/gitflow-workflow}
            \item \url{https://github.com/about}
            \item \url{https://www.atlassian.com/software/jira}
            \item \url{https://www.atlassian.com/}
            \item \url{https://meet.google.com/}
            \item \url{https://slack.com/intl/en-it/about}
            \item \url{https://telegram.org/}
            \item \url{https://workspace.google.com/products/chat/}
        \end{itemize}

\pagebreak

	\pagebreak

	\section{Tecnologie e librerie}

Nelle sezioni che seguono, vengono elencate le tecnologie e le librerie interessate dallo sviluppo del software. Per ognuna, viene presentata una breve descrizione e spiegato il suo impiego nel contesto del software. Dove necessario, viene fornito un collegamento per il download e l'installazione delle risorse necessarie per lo sviluppo e manutenzione del progetto.





















	\pagebreak

	
	\section{Setup}
PORTACS viene distribuito tramite container Docker, per cui i dispositivi sui quali dovrà eseguire avranno minimi requisiti software. È possibile utilizzare i container anche per la fase\textsubscript{G} di sviluppo, altrimenti si possono scaricare ed utilizzare gli strumenti descritti in \ref{tecnologie} per l'esecuzione diretta in locale. PORTACS\textsubscript{A} si divide su tre immagini Docker:
\begin{enumerate}
    \item server;
    \item client muletto (forklift);
    \item client utente (user).
\end{enumerate}


\subsection{Requisiti di sistema}
Sotto elencati saranno descritti i requisiti minimi del sistema per un corretto funzionamento del software PORTACS\textsubscript{A}.


\subsubsection{Requisiti Hardware}
\begin{itemize}
	\item Client → unità:
\begin{itemize}
	\item CPU\textsubscript{A} dual-core o maggiore con frequenza indicativa >= 3 Ghz;
	\item memoria RAM\textsubscript{A} >= 4GB dedicati;
	\item connessione con bassi tempi di risposta.
\end{itemize}
	\item Client → admin o responsabile:
\begin{itemize}
	\item CPU\textsubscript{A} dual-core o maggiore con frequenza indicativa >= 3 Ghz;
	\item memoria RAM\textsubscript{A} >= 4GB dedicati;
	\item connessione con bassi tempi di risposta.
\end{itemize}
	\item Server:
\begin{itemize}
	\item CPU\textsubscript{A} quad-core o maggiore con frequenza indicativa >= 4 Ghz;
	\item memoria RAM\textsubscript{A} >= 8GB dedicati;
	\item connessione con bassi tempi di risposta.
\end{itemize}
\end{itemize}

\subsubsection{Requisiti Software}
	\begin{comment}
        \pparagraph{Librerie}
        \begin{itemize}
        \item express (v4.17.1);
        \item http (v0.0.1-security);
        \item mocha (v8.3.2);
        \item net (v1.0.2);
        \item primeicons (v4.1.0);
        \item primeng (v11.3.2);
        \item rxjs (v6.6.0);
        \item socket.io (v3.1.2);
        \item spring boot (v2.4.5).
        \end{itemize}
    \end{comment}

    \pparagraph{Esecuzione}
    \begin{itemize}
   		\item Utilizzando Docker, il sistema operativo host non dovrebbe essere un problema, tuttavia in seguito si daranno le istruzioni per l'esecuzione valide in ambiente unix, nel quale si è anche effettuato la maggior parte dei test (Linux kernel v5.8 o MacOS v10.15);
        \item Docker (v19.03.*), installazione base, per esempio in Linux è sufficiente: \texttt{sudo snap install docker};
        \item Google Chrome (v90).
    \end{itemize}

    \pparagraph{Sviluppo}
    Vedi \S\ \ref{tecnologie}

    \begin{comment}
    	\item Client:
    \begin{itemize}
    	\item Docker.
    	%\item Chrome Versione 90.0.4430.85 (Build ufficiale) (a 64 bit);
    	%\item Node.js;
    	%\item Angular;

    	%\item Windows 10.
    \end{itemize}
    	\item Server:
    \begin{itemize}
    	\item Docker.
    	%\item Java;
    	%\item gradle;

    	%\item Windows 10.
    \end{itemize}
    \end{comment}







\subsection{Installazione ed avvio degli applicativi}
    \subsubsection{Server}
    Nella macchina da utilizzare come server andrà scaricata l'immagine di portacs-server con la versione desiderata dal \href{https://hub.docker.com/r/threewaymilkshake}{docker hub di \group}.
    Per l'avvio, in ambiente Linux e MacOS:
    \begin{itemize}
        \item predisporre una cartella nella quale posizionarsi per l'esecuzione che contenga una cartella \texttt{database} che manterrà la persistenza, questa dovrà contenere i file \texttt{map.json}, \texttt{forklifts.json} e \texttt{users.json};
        \item posizionarsi su tale cartella;
        \item eseguire:
    \begin{verbatim}
    docker run --rm --name server \
        -v $(pwd)/database:/database \
        -p 1723:1723 \
        threewaymilkshake/portacs-server
    \end{verbatim}
    \end{itemize}
    Il container verrà avviato nella finestra corrente e si potranno osservare i log del server, questo può essere terminato dalla stessa finestra con la combinazione \texttt{Ctrl+C}. Se si vuole eseguire il container in background è sufficiente aggiungere l'opzione \texttt{-d}, per terminarlo a questo punto bisognerà eseguire: \texttt{docker stop server}.
    Si raccomanda di impostare un ip statico per la macchina server, così da non dover richiedere variazioni frequenti nelle configurazioni dei client.

    \subsubsection{Client}
    Come per il server, scaricare l'immagine dal \href{https://hub.docker.com/r/threewaymilkshake}{docker hub di \group{}} relativa al client voluto (forklift o user). Dopodiché posizionarsi in una cartella dove predisporre un file di configurazione \texttt{config.txt}:
    \begin{itemize}
        \item per i client \textit{user} sarà sufficiente specificare l'indirizzo IP della macchina server così: \texttt{SERVER\_ADDR=ip};
        \item per i client \texttt{forklift}, oltre all'indirizzo del server come per gli utenti, bisognerà aggiungere altre informazioni:
        \begin{itemize}
            \item \textbf{obbligatorie: }
            \begin{itemize}
                \item id e token del muletto;
                \item posizione iniziale (corrispondendte alla base dove il muletto verrà posizionato al primo avvio);
            \end{itemize}
            \item \textbf{opzionali: }
            \begin{itemize}
                \item per specificare delle porte diverse da far utilizzare ai servizi, quest'operazione è necessaria solo se si vogliono simulare più unità nella stessa macchina locale, al fine di evitare collisioni tra le porte;
                \item bisogna impostare due campi, \texttt{ngPort1} e \texttt{ngPort2}, devono essere diversi da altre configurazioni in uso sulla stessa macchina;
                \item bisogna predisporre un file diverso per ogni unità che si vuole simulare.
            \end{itemize}
        \end{itemize}
        Dev'essere poi sempre presente \texttt{configurationEnv=custom}. L'ordine in cui vengono specificati i campi non è importante. Segue un esempio di configurazione ben formata con utilizzo delle impostazioni opzionali:
        \begin{verbatim}
            startX=0
            startY=1
            ngPort1=4203
            ngPort2=8083
            forkliftId=f3
            forkliftToken=abc
            configurationEnv=custom
            SERVER_ADDR=192.168.3.129
        \end{verbatim}
        Queste informazioni sono a disposizione degli admin.
    \end{itemize}

    Per avviare un muletto sarà sufficiente eseguire (i valori di default per \texttt{ngPort1} e \texttt{ngPort2} sono rispettivamente \texttt{4201} e \texttt{8081}):
    \begin{verbatim}
        docker run --rm --name <nome> -d \
            --env-file config.txt
            -p 4201:4201 \
            -p 8081:8081 \
            portacs-client-forklift
    \end{verbatim}
    Il <nome> specificato dopo \texttt{---name} servirà per terminare il container.

    Per avviare un utente invece:
    \begin{verbatim}
        docker run --rm --name user -d \
            --env-file user.txt \
            -p 4300:4300 \
            -p 8090:8090 \
            portacs-client-user
    \end{verbatim}

    \pparagraph{Terminazione dei container}
        Se si vuole terminare un container l'istruzione è: \texttt{docker stop <nome>}, dove <nome> è il parametro che si ha specificato dopo \texttt{---name}. La specifica del nome all'avvio del container non è obbligatoria, nel caso di esecuzione con opzione \texttt{-d} verrà stampato l'id che si potrà usare al posto del nome per terminare il container o in ogni caso si possono ispezionare i container attivi tramite \texttt{docker ps}.

\clearpage
\subsection{Costruzione degli artefatti e delle immagini}
    Quando il codice sorgente viene modificato le immagini docker corrispondenti devono essere ricostruite e prima di ciò bisogna produrre gli artefatti necessari. Per le operazioni descritte in seguito si assume di avere le cartelle con i sorgenti relativi.
    \subsubsection{Server}
        Per produrre l'unico artefatto\textsubscript{G} jar necessario, si deve eseguire \texttt{./gradlew build}, questo produrrà appunto l'eseguibile segnandolo con la versione specificata in \texttt{build.gradle}. Questo comando esegue anche tutti i test di unità segnalando eventuali fallimenti. La build ha come precondizione\textsubscript{G} la corretta formattazione del codice, la quale può essere ottenuta con \texttt{./gradlew spotlessApply}.
        Con il jar aggiornato si può eseguire la build di docker per ottenere l'immagine aggiornata:
        \begin{verbatim}
            docker build -t threewaymilkshake/portacs-server .
        \end{verbatim}

    \subsubsection{Client}
        Per entrambi i client bisogna assicurarsi che il file \texttt{node\_package.json} nelle cartelle corrispondenti contenga tutte e sole le dipendenze necessarie a node. Dopodiché bisogna compilare la parte angular:
        \begin{itemize}
            \item per i muletti: \texttt{ng build -c dev1};
            \item per gli utenti: \texttt{ng build}.
        \end{itemize}
        A questo punto si possono generare le rispettive immagini:
        \begin{verbatim}
            docker build -t threewaymilkshake/portacs-client-forklift .
            docker build -t threewaymilkshake/portacs-client-user .
        \end{verbatim}


    \subsubsection{Caricamento nel docker hub}
        \noindent Per aggiornare le immagini nel docker hub, è necessario avere le credenziali o essere collaboratori. Dopo aver effettuato l'accesso tramite tramite \texttt{docker login}, è sufficiente eseguire: \texttt{docker {push} <nome-immagine>}













	\pagebreak
	

	\section{Testing}
Questo capitolo ha lo scopo di indicare agli sviluppatori come controllare in che modo opera il codice e la sua sintassi. Vengono di seguito esposti gli strumenti utilizzati per effettuare i test di analisi statica e dinamica del nostro applicativo.

\subsection{Coveralls}
Per quanto riguarda l'attività di code coverage è stato scelto di utilizzare \textit{Coveralls}. \textit{Coveralls} esegue revisioni automatiche al codice relative all'analisi statica attraverso \textit{GitHub Actions}.

\subsection{GitHub Actions}
Il servizio di Continuous Integration che è stato deciso di utilizzare è \textit{GitHub Actions}, fornito da \textit{GitHub}. \\\textit{GitHub Actions} permette di creare dei workflow, ovvero processi automatici, con l'obiettivo di automatizzare il ciclo di vita dello sviluppo del software.

% CI e code coverage, glossario




	\newpage


	\section{Architettura del sistema}


Qui si potrebbe mettere uno schema simile a quello della slide iniziale per evidenziare l'architettura client-server



\subsection{Server}

Dire che è 3 layer architecture

Qui potrebbe esserci il diagramma minimale complessivo


\subsubsection{Diagramma delle classi}


\paragraph{Persistance layer}


\paragraph{Business layer}

\subparagraph{Mappa}

\subparagraph{Clients}

\subparagraph{Tasks}

\subparagraph{Collisioni}


\paragraph{Communication layer}



\subsection{Client}


\subsection{Comunicazione}

\subsubsection{Diagrammi di sequenza}

\subsubsection{Protocollo di comunicazione}



	\newpage


	\section{Estendere PORTACS}
In questa sezione verranno riportate tutte quelle informazioni utili ad una semplice e corretta estensione del prodotto PORTACS. Possono essere estensioni legate sia a nuovi algoritmi e framework, più efficaci o efficienti, che a nuove categorie di sotto-sistemi di PORTACS.















	\newpage


\end{document}
